\documentclass[11pt]{article}
\usepackage{amsmath,amssymb,amsthm}
\usepackage[margin=1in]{geometry}
\usepackage{array}
\usepackage{tabularx}
\usepackage{booktabs}
\usepackage[hidelinks]{hyperref}

\theoremstyle{plain}
\newtheorem{theorem}{Theorem}[section]
\newtheorem{lemma}[theorem]{Lemma}
\newtheorem{proposition}[theorem]{Proposition}
\newtheorem{corollary}[theorem]{Corollary}
\theoremstyle{definition}
\newtheorem{definition}[theorem]{Definition}
\newtheorem{example}[theorem]{Example}
\newtheorem{fact}[theorem]{Fact}
\theoremstyle{remark}
\newtheorem{remark}[theorem]{Remark}

\newcommand{\Z}{\mathbb{Z}}
\newcommand{\sm}{\setminus}
\DeclareMathOperator{\exc}{exc}
\DeclareMathOperator{\solo}{solo}
\newcommand{\Adm}{\mathrm{Adm}}
\newcommand{\dg}{\operatorname{deg}}

\newcommand{\partheading}[1]{%
  \par\vspace{1.6\baselineskip}%
  \noindent\rule{\linewidth}{0.8pt}\par\vspace{0.35\baselineskip}%
  \begingroup\centering\normalsize\scshape #1\par\endgroup
  \vspace{0.35\baselineskip}\noindent\rule{\linewidth}{0.8pt}%
  \par\vspace{1.0\baselineskip}}

\title{A one-cell deletion proof for bounded-excess growth of minimizing polyhexes}
\author{H. Utsunomiya}
\date{}
\begin{document}
\maketitle
\begin{abstract}
This note records a second proof of the one-cell growth theorem of the author's paper
\emph{Shortest-side skeletons of edge-perimeter-minimizing polyhexes}. Instead of deleting
a whole extreme line at a time, it deletes one cell at a time, keeping the configuration
directionally convex with at most one occupied line above the threshold and tracking two
competing budgets through an admissibility condition. The existence of a legal deletion and
the preservation of the condition are proved separately; the former rests on a
bridge-excess lemma, the latter on a bound at the start of each plateau of the threshold,
whose two exceptional cases are the rigidity statement used by the main paper. Two
configurations show that neither auxiliary condition can be dropped. The common framework
of the main paper is reproduced for self-containment.
\end{abstract}

The related construction of Brunvoll, Cyvin and Cyvin
\cite[Theorems~1--2]{BrunvollCyvinCyvin1993} uses
\emph{circumscribing} (adding a closed peripheral ring of cells) and its
inverse, \emph{excising}, to study eventual extremality when repeated
circumscribing is possible. These whole-ring operations differ from
the single-cell deletion with a bounded-excess invariant proved here.

\section{The hexagonal lattice, runs, and directional convexity}\label{sec:prelim}

This part reproduces, for self-containment, the common framework of the main paper
[Utsunomiya, \emph{Shortest-side skeletons of edge-perimeter-minimizing polyhexes}]:
the lattice, the perimeter functional, directional convexity, the tight hull and its
identities, and the extremal capacity function. Numbering follows this note.

\subsection{Lattice and polyhexes}\label{sec:2.1}

Cells of the hexagonal tiling are indexed by axial coordinates
\((q,r)\in\Z^{2}\). We call \(q\), \(r\) and
\[
   s:=q+r
\]
the three \emph{keys} of a cell; Two cells are \emph{adjacent} when they differ by one of the six
vectors
\[
   d_{0}=(1,0),\quad d_{1}=(0,1),\quad d_{2}=(-1,1),\quad
   d_{3}=(-1,0),\quad d_{4}=(0,-1),\quad d_{5}=(1,-1).
\]

A \emph{polyhex} is a finite nonempty set of cells that is connected with respect to
adjacency. We write \(|S|\) for the number of cells of \(S\).

The following four facts are immediate from the definitions and are used throughout
without further comment.

\begin{fact}\label{f:1}
Two distinct cells are adjacent if and only if they agree in exactly one key and differ
by exactly \(1\) in the other two. In words: an adjacency step keeps one of the three
keys fixed and changes each of the other two by one unit (the step \(d_{0}=(1,0)\), for
instance, sends \((q,r,s)\) to \((q+1,r,s+1)\), so the two changing keys need not move
in opposite senses).
\end{fact}

\begin{proof}
Each \(d_{j}\) changes \((q,r,s)\) by \((1,0,1)\), \((0,1,1)\), \((-1,1,0)\),
\((-1,0,-1)\), \((0,-1,-1)\) or \((1,-1,0)\), and these are exactly the six triples with
one zero entry and two entries equal to \(\pm1\) that are realized. Conversely a
difference of that shape is one of the \(d_{j}\).
\end{proof}

\begin{fact}\label{f:2}
A \emph{line} is a set of cells on which one key is constant. Two lines belonging to
different keys meet in exactly one cell; two distinct lines belonging to the same key are
disjoint.
\end{fact}

\begin{fact}\label{f:3}
One adjacency step changes each of \(q\), \(r\) and \(s\) by at most \(1\).
\end{fact}

\begin{fact}\label{f:4}
For a polyhex \(S\) and each key, the set of values taken by that key on \(S\) is an
interval of consecutive integers. Thus every intermediate key value occurs in \(S\).
\end{fact}

\begin{proof}
Let \(k\) be a key and let \(x,y\in S\) realize two values of \(k\). A path of adjacent
cells inside \(S\) joins them, and by Fact~\ref{f:3} the value of \(k\) changes by at most
one per step, so every intermediate integer occurs on the path.
\end{proof}

\subsection{Edge perimeter, runs and directional convexity}\label{sec:2.2}

Let \(a(S)\) denote the number of unordered pairs of adjacent cells of \(S\). The
\emph{edge perimeter} of \(S\) is
\begin{equation}\label{eq:p}
   p(S)=6|S|-2a(S),
\end{equation}
Equivalently, \(p(S)\) counts the cell edges of \(S\) that are not shared with another
cell of \(S\); we call these the \emph{boundary edges}. The \emph{edge semiperimeter} is
\[
   \sigma(K):=\tfrac{1}{2}p(K)
\]
for any finite set of cells \(K\). Its integrality follows from Lemma~\ref{lem:runs}.

A line \(\Lambda\) is \emph{occupied} by \(S\) when \(S\cap\Lambda\neq\emptyset\); we call
\(S\cap\Lambda\) the \emph{trace} of \(S\) on \(\Lambda\). Cells of a line are linearly
ordered (for a line of constant \(q\) by the value of \(r\), for one of constant \(r\) by
the value of \(q\), for one of constant \(s\) by the value of \(q\)), and consecutive
cells in that order are adjacent. A \emph{run} of \(S\) on \(\Lambda\) is a maximal block
of consecutive cells of the trace. Let
\[
   R(S)=\text{the total number of runs of \(S\) over all lines},\qquad
   \lambda(S)=\text{the number of occupied lines}.
\]

\begin{definition}\label{def:dconv}
A polyhex \(S\) is \emph{directionally convex} when the trace of \(S\) on every occupied
line is a single run. Equivalently, each occupied line meets \(S\) in an interval of cells.
\end{definition}

This is a lattice notion and not Euclidean convexity; it is the notion called
\emph{layer convexity} in parts of the literature on lattice animals.

\begin{lemma}\label{lem:runs}
For every finite set of cells \(S\),
\begin{equation}\label{eq:pruns}
   p(S)=2R(S)\ \ge\ 2\lambda(S),
\end{equation}
In particular \(p(S)\) is even and
\(\sigma(S)=R(S)\). For a polyhex, equality \(R(S)=\lambda(S)\) holds if and only if
\(S\) is directionally convex.
\end{lemma}

\begin{proof}
Let \(e=\{x,y\}\) be a boundary edge, with \(x\in S\) and \(y\notin S\) the two cells it
separates. By Fact~\ref{f:1} the cells \(x\) and \(y\) agree in exactly one key, hence
they are consecutive cells of exactly one line \(\Lambda\), and \(x\) is an end of the run
of \(S\) on \(\Lambda\) containing it. Conversely a run has exactly two ends, and the cell
immediately beyond each end lies outside \(S\) and determines a boundary edge; for a run
consisting of a single cell the two ends are the two opposite sides of that cell, so they
give two distinct boundary edges. This is a bijection between boundary edges and ends of
runs, whence \(p(S)=2R(S)\).

Every occupied line carries at least one run, so \(R(S)\ge\lambda(S)\), with equality
exactly when every occupied line carries exactly one run, which is
Definition~\ref{def:dconv}.
\end{proof}
\noindent The run-count identity and its slice-convex equality case were
previously formulated by Yang \cite[Theorems~4--5]{Yang2002}.


\begin{remark}[edge perimeter versus site perimeter]\label{rem:site}
The \emph{site perimeter} of \(S\) is the number of cells outside \(S\) adjacent to at
least one cell of \(S\). It is a different functional from \eqref{eq:p}, and the two have
different minimizers in general; results about one do not transfer to the other. Nothing
in this paper depends on the site perimeter. We mention it only because part of the
literature we cite in Section~\ref{sec:2.5} is written in those terms.
\end{remark}

\subsection{Extreme support lines and traces}\label{sec:2.3}

Let \(S\) be a polyhex. For each key, Fact~\ref{f:4} makes the occupied values an interval,
so there is a smallest and a largest occupied line in each of the three families. These six
lines are the \emph{extreme lines} of \(S\); we list them in the cyclic order
\[
   \Lambda_{1}:\ q=q_{\min},\quad
   \Lambda_{2}:\ r=r_{\max},\quad
   \Lambda_{3}:\ s=s_{\max},\quad
   \Lambda_{4}:\ q=q_{\max},\quad
   \Lambda_{5}:\ r=r_{\min},\quad
   \Lambda_{6}:\ s=s_{\min},
\]
so that \(\Lambda_{i}\) and \(\Lambda_{i+3}\) are the two extreme lines of one and the same
family, indices being read in \(\Z/6\). If a family occupies a single line then the two
corresponding extreme lines coincide; all statements below are read with that degeneracy
allowed.

\begin{definition}\label{def:traces}
Let \(S\) be a \emph{directionally convex} polyhex. Its six \emph{support traces} are the
traces \(S\cap\Lambda_{i}\), each a single run by Definition~\ref{def:dconv}, and we write
\[
   r_{i}:=|S\cap\Lambda_{i}|\ \ge 1 ,\qquad i\in\Z/6,
   \qquad\text{and}\qquad
   \rho(S):=\min_{i} r_{i} .
\]
Thus \(\rho(S)\) is the minimum support-trace length.
\end{definition}

\begin{remark}\label{rem:rho}
Definition~\ref{def:traces} is stated for directionally convex \(S\) only. On that domain \(\rho(S)\) is at once the length of the shortest support trace and the
length of the shortest extreme-line run. For
a general polyhex a trace need not be a run, and we do not extend the definition.
\end{remark}

\section{Tight hulls and exact identities}\label{sec:hulls}

\subsection{Tight lattice hulls}\label{sec:2.4}

\begin{definition}\label{def:hexagon}
Let \(L\ge1\) and \(u,v,w\ge0\) be integers satisfying the \emph{admissibility} conditions
\begin{equation}\label{eq:adm}
   u+v\le L-1,\qquad v+w\le L-1,\qquad u+w\le L-1 ,
\end{equation}
and put \(A:=L-u\), \(B:=L-v\), \(C:=L-w\), so that \(A,B,C\ge1\). The corresponding
\emph{lattice hexagon} is
\[
   H(L;u,v,w):=\{\,(q,r):\ 0\le q\le A-1,\ \ 0\le r\le B-1,\ \ w\le q+r\le L-1\,\}.
\]
Geometrically, this is a triangle of side \(L\) with corner triangles of sides \(u\), \(v\), and
\(w\) removed. Degenerate shapes, in which one or more of the cut corners is empty or one
of the three widths equals \(1\), are allowed.
\end{definition}

\begin{definition}\label{def:hull}
The \emph{tight hull} \(H(S)\) of a polyhex \(S\) is the lattice hexagon determined by the
extreme values of the three keys on \(S\). Concretely, normalize by a translation so that
\(q_{\min}=r_{\min}=0\) and set
\[
   A=1+q_{\max},\quad B=1+r_{\max},\quad L=1+s_{\max},\quad w=s_{\min},\quad
   u=L-A,\quad v=L-B .
\]
\end{definition}

\begin{lemma}[the tight hull is admissible]\label{lem:tightadm}
For every polyhex \(S\), the parameters \(L,u,v,w\) of Definition~\ref{def:hull} satisfy
\(L\ge1\), \(u,v,w\ge0\) and the three inequalities \eqref{eq:adm}, so
\(H(S)=H(L;u,v,w)\) is a lattice hexagon.
\end{lemma}

\begin{proof}
After the normalization every cell has \(q,r\ge0\), hence \(s\ge0\), so \(w\ge0\) and
\(L\ge1\); and \(u=s_{\max}-q_{\max}\), \(v=s_{\max}-r_{\max}\). A cell with
\(q=q_{\max}\) has \(s\ge q_{\max}\), so \(u\ge0\); likewise \(v\ge0\). A cell with
\(s=s_{\max}\) has \(s_{\max}=q+r\le q_{\max}+r_{\max}\), which is \(u+v\le L-1\). A cell
with \(q=0\) has \(s=r\le r_{\max}\), so \(s_{\min}\le r_{\max}\), which is \(v+w\le L-1\);
a cell with \(r=0\) gives \(s_{\min}\le q_{\max}\), which is \(u+w\le L-1\).
\end{proof}

For a lattice hexagon \(H\) we write \(\Lambda_{1},\dots,\Lambda_{6}\) for its six extreme
lines in the cyclic order of Section~\ref{sec:2.3}, namely
\[
   q=0,\quad r=B-1,\quad s=L-1,\quad q=A-1,\quad r=0,\quad s=w ,
\]
and call them the \emph{side lines}; we write
\[
   \ell_{i}:=|H\cap\Lambda_{i}|
\]
for the six \emph{side lengths}, \(h:=\min_{i}\ell_{i}\) for the shortest side length, and
\(\kappa_{i}:=H\cap\Lambda_{i}\cap\Lambda_{i+1}\) for the \emph{corner} cells. Finally we
set
\[
   \delta(H):=\ell_{1}-\ell_{4},
   \qquad
   Q(H):=(\ell_{1}-\ell_{3})^{2}+(\ell_{3}-\ell_{5})^{2}+(\ell_{5}-\ell_{1})^{2} .
\]

\begin{lemma}\label{lem:hull}
Let \(S\) be a polyhex and \(H=H(S)\). Then
\begin{enumerate}
\item[(i)] \(S\subseteq H\);
\item[(ii)] \(S\) and \(H\) occupy exactly the same lines, so \(\lambda(S)=\lambda(H)\);
\item[(iii)] each of the six side lines of \(H\) meets \(S\).
\end{enumerate}
\end{lemma}

\begin{proof}
(i) is immediate from Definition~\ref{def:hull}, since every cell of \(S\) satisfies the
six inequalities defining \(H\). For (ii), the lines occupied by \(S\) in the \(q\)-family
are, by Fact~\ref{f:4}, exactly \(q=0,\dots,A-1\), and likewise \(r=0,\dots,B-1\) and
\(s=w,\dots,L-1\); these are precisely the lines occupied by \(H\), as read off from
Definition~\ref{def:hexagon}. Assertion (iii) is the case of (ii) for the six extreme
values: each of them is attained by some cell of \(S\).
\end{proof}

Lemma~\ref{lem:hull} is the only place where connectedness of \(S\) is used in this
section, and it is used through Fact~\ref{f:4}: for a disconnected set the occupied key
values need not form an interval, and (ii) fails.

\subsection{Hull identities}\label{sec:2.5}

\begin{lemma}\label{lem:hex}
Let \(H=H(L;u,v,w)\) be a lattice hexagon. Then:
\begin{enumerate}
\item[(i)] the six side lengths are
\[
   \ell_{1}=L-v-w,\quad \ell_{2}=v+1,\quad \ell_{3}=L-u-v,\quad
   \ell_{4}=u+1,\quad \ell_{5}=L-u-w,\quad \ell_{6}=w+1 ,
\]
and each side line meets \(H\) in a single run whose two end cells are the corners
\(\kappa_{i-1}\) and \(\kappa_{i}\);
\item[(ii)] \(H\) is a directionally convex polyhex with
\[
   \lambda(H)=\sigma(H)=A+B+C=3L-(u+v+w);
\]
\item[(iii)] the six side lengths satisfy
\begin{equation}\label{eq:sidesum}
   \sum_{i=1}^{6}\ell_{i}=\lambda(H)+3 ,
\end{equation}
in words: the six side lengths of the hull add up to three more than the number of lines
the hull occupies; and the three differences across opposite sides agree,
\begin{equation}\label{eq:delta}
   \ell_{1}-\ell_{4}=\ell_{3}-\ell_{6}=\ell_{5}-\ell_{2}=\delta(H)=L-u-v-w-1 ;
\end{equation}
\item[(iv)] with \(T(k):=k(k+1)/2\),
\[
   |H|=T(L)-T(u)-T(v)-T(w) ;
\]
\item[(v)] \textup{(hull defect identity)}
\begin{equation}\label{eq:defectid}
   \lambda(H)^{2}+3-12|H| \;=\; 3\,\delta(H)^{2}+2\,Q(H)\;\ge\;0 ,
\end{equation}
in words: the square of the number of occupied lines, plus three, minus twelve times the
number of cells, equals three times the square of the common opposite-side difference plus
twice the imbalance \(Q(H)\).
\end{enumerate}
\end{lemma}

\begin{proof}
(i) On the line \(q=0\) the defining inequalities reduce to \(w\le r\le\min(B-1,L-1)=B-1\),
which is a nonempty interval because \(w\le B-1\) is the condition \(v+w\le L-1\) of
\eqref{eq:adm}; its length is \(B-w=L-v-w\). On the line \(r=B-1\) they reduce to
\(\max(0,w-B+1)\le q\le\min(A-1,v)\), and \eqref{eq:adm} gives \(w-B+1\le0\) and
\(v\le A-1\), so the interval is \(0\le q\le v\), of length \(v+1\). The remaining four
sides are obtained in the same way, each time using exactly one of the three inequalities
of \eqref{eq:adm} at each end. Each side is an interval of consecutive cells of its line,
hence a run; its two end cells are the two cells where a second defining inequality becomes
an equality, that is, the two corners.

(ii) The occupied lines are \(q=0,\dots,A-1\), \(r=0,\dots,B-1\) and \(s=w,\dots,L-1\),
which number \(A+B+C\); on each of them the trace is an interval of consecutive cells by
the same computation as in (i), so \(H\) is directionally convex and
Lemma~\ref{lem:runs} gives \(\sigma(H)=R(H)=\lambda(H)\). Connectedness follows since each
trace is a run and consecutive traces meet a common line.

(iii) Summing the six expressions of (i),
\[
   \sum_{i}\ell_{i}=3L-(u+v+w)+3=\lambda(H)+3 .
\]
Using the side-length formulas of (i), the three opposite-side differences
\(\ell_{1}-\ell_{4}\), \(\ell_{3}-\ell_{6}\), \(\ell_{5}-\ell_{2}\) of \eqref{eq:delta}
are \((L-v-w)-(u+1)\), \((L-u-v)-(w+1)\) and \((L-u-w)-(v+1)\), all equal to
\(L-u-v-w-1\).

(iv) The triangle \(\{q,r\ge0,\ q+r\le L-1\}\) has \(T(L)\) cells. The three sets removed
in Definition~\ref{def:hexagon} are the corner triangles \(q\ge A\), \(r\ge B\) and
\(q+r\le w-1\), of \(T(u)\), \(T(v)\) and \(T(w)\) cells respectively, and \eqref{eq:adm}
makes them pairwise disjoint: a cell with \(q\ge A\) and \(r\ge B\) would have
\(q+r\ge A+B=2L-u-v\ge L+1\), which is impossible, and a cell with \(q\ge A\) and
\(q+r\le w-1\) would give \(r\le w-1-A<0\), and symmetrically.

(v) Put \(U:=u+v+w\) and \(P:=u^{2}+v^{2}+w^{2}\). By (ii) and (iv),
\[
   \lambda(H)^{2}+3-12|H|=(3L-U)^{2}+3-6L(L+1)+6P+6U
   =3L^{2}-6LU-6L+U^{2}+6P+6U+3 .
\]
On the other side, \(\delta(H)=L-U-1\) and, by (i),
\(Q(H)=(u-w)^{2}+(w-v)^{2}+(v-u)^{2}=3P-U^{2}\), so
\[
   3\delta(H)^{2}+2Q(H)=3(L-U-1)^{2}+2(3P-U^{2})
\]
expands to the same expression. In words: both sides equal three times \(L^{2}\), minus six
times \(LU\), minus six times \(L\), plus \(U^{2}\), plus six times \(P\), plus six times
\(U\), plus three. Nonnegativity is clear from the right-hand side.
\end{proof}



\begin{remark}[attribution]\label{rem:fs}
An equivalent defect identity is obtained from F\"ulep and Sieben's polyiamond identity
\cite[Lemma~3.3]{FulepSieben2010} through the triangular--hexagonal duality developed in their Section~5.
In the same site-perimeter variables, \cite[Proposition~5.10]{FulepSieben2010} yields exactly the capacity
function \(M(t)\) of Section~\ref{sec:2.6} after the substitution \(p=t+3\); the value
\(p=7\) excluded there corresponds to the value \(t=4\) excluded in
Remark~\ref{rem:attain} below. We nevertheless give the derivation here in edge-perimeter
variables to keep the argument self-contained, and we do not identify the two settings:
\cite{FulepSieben2010} is about the site perimeter (Remark~\ref{rem:site}).
\end{remark}

\section{Capacity and minimum edge perimeter}\label{sec:capacity}

\subsection{Extremal capacity}\label{sec:2.6}

We first record the two operations on directionally convex polyhexes that the capacity
computation needs: removing a whole support trace, and putting it back one cell at a time.

\begin{lemma}[peeling]\label{lem:peel}
Let \(K\) be a directionally convex polyhex, let \(\Lambda\) be one of its six extreme
lines, and let \(K':=K\sm\Lambda\). If \(K'\neq\emptyset\), then
\begin{enumerate}
\item[(i)] \(K'\) is a directionally convex polyhex with \(\sigma(K')\le\sigma(K)-1\);
\item[(ii)] the cells of \(K\cap\Lambda\) can be ordered \(c_{1},\dots,c_{k}\) so that for
every \(j\) the set \(K'\cup\{c_{1},\dots,c_{j}\}\) is a directionally convex polyhex of
edge semiperimeter at most \(\sigma(K)\).
\end{enumerate}
\end{lemma}

\begin{proof}
Say \(\Lambda\) is the line \(q=q_{\min}\); the other five cases are the same after
permuting the roles of the keys and of the two senses.

(i) For each occupied value \(k\) of \(q\) let \(K_{k}\) be the trace of \(K\) on
\(q=k\); it is a run. If \(K\) is connected, then for consecutive values \(k,k+1\) a path
inside \(K\) joining a cell of \(K_{k}\) to a cell of \(K_{k+1}\) must, by
Fact~\ref{f:3}, contain a step from a cell with \(q=k\) to a cell with \(q=k+1\), so
\(K_{k}\) and \(K_{k+1}\) are adjacent. Hence \(K'\), the union of the \(K_{k}\) with
\(k\ge q_{\min}+1\), is connected. Its traces in the \(q\)-family are traces of \(K\),
hence runs. Now fix a line of the \(r\)-family, \(r=k\); its cells are \((q,k)\), so the
trace of \(K\) on it is a run if and only if its \(q\)-values form an interval, and by
Fact~\ref{f:2} it meets \(\Lambda\) in at most the one cell \((q_{\min},k)\). The trace of
\(K'\) on this line is the trace of \(K\) with that cell, if present, removed. Since
\(q_{\min}\) is the least \(q\)-value occurring anywhere in \(K\), it is the least
\(q\)-value of the interval, so removing it leaves an interval, possibly empty. A line of
the \(s\)-family, \(s=k\), has cells \((q,k-q)\), again parametrized by \(q\), and the
same sentence applies verbatim. So all traces of \(K'\) are runs and \(K'\) is
directionally convex. Finally \(\Lambda\) is occupied by
\(K\) and not by \(K'\), so \(\lambda(K')\le\lambda(K)-1\), and by Lemma~\ref{lem:runs}
\(\sigma(K')=\lambda(K')\le\lambda(K)-1=\sigma(K)-1\).

(ii) Since \(K\) is connected there is a cell \(c_{1}\in K\cap\Lambda\) adjacent to
\(K'\). Order the remaining cells of the run \(K\cap\Lambda\) so that every initial
segment \(\{c_{1},\dots,c_{j}\}\) is a block of consecutive cells of that run containing
\(c_{1}\). Write \(K_{j}:=K'\cup\{c_{1},\dots,c_{j}\}\). Then \(K_{j}\) is connected,
because \(c_{1}\) is adjacent to \(K'\) and each later \(c_{j}\) is adjacent to a cell
already inserted. Its trace on \(\Lambda\) is a block, hence a run; its trace on any other
line of the \(q\)-family is a trace of \(K'\), hence a run. On a line of the \(r\)- or
\(s\)-family, parametrized by \(q\) as in (i), the trace of \(K_{j}\) is the trace of
\(K'\) there together with at most one inserted cell, which lies on \(\Lambda\) and so
has \(q=q_{\min}\). The trace of \(K\) on that line is a run containing this cell and, by
(i), the trace of \(K'\) is that run with its least element \(q_{\min}\) removed; putting
the cell back restores the run of \(K\). Hence every trace of \(K_{j}\) is a run, so
\(K_{j}\) is directionally convex, and every line it occupies is occupied by \(K\), whence
\(\sigma(K_{j})=\lambda(K_{j})\le\lambda(K)=\sigma(K)\).
\end{proof}

\begin{corollary}\label{cor:subsets}
Let \(K\) be a directionally convex polyhex. For every \(j\) with \(1\le j\le|K|\) there is
a directionally convex polyhex \(K_{j}\subseteq K\) with \(|K_{j}|=j\) and
\(\sigma(K_{j})\le\sigma(K)\).
\end{corollary}

\begin{proof}
Strong induction on \(|K|\). For \(|K|=1\) there is nothing to prove. Let \(|K|\ge2\). Some
extreme line \(\Lambda\) does not contain all of \(K\): otherwise \(K\) would be contained
in two extreme lines belonging to different families, which by Fact~\ref{f:2} meet in a
single cell, contradicting \(|K|\ge2\). Apply Lemma~\ref{lem:peel} to that \(\Lambda\). The
sets \(K_{j}\) of part (ii) supply all sizes between \(|K'|\) and \(|K|\); the sizes below
\(|K'|\) are supplied by the inductive hypothesis applied to \(K'\), whose edge
semiperimeter is at most \(\sigma(K)\).
\end{proof}

\begin{definition}\label{def:M}
For an integer \(t\ge0\) put
\[
   M(t):=\left\lfloor\frac{t^{2}+3}{12}\right\rfloor ,
\]
Thus
\(M(0)=M(1)=M(2)=0\) and \(M(3)=1\). The function \(M\) is nondecreasing.
\end{definition}

\begin{theorem}\label{thm:capacity}
Let \(N\ge1\).
\begin{enumerate}
\item[(a)] Every polyhex \(S\) satisfies \(|S|\le M(\sigma(S))\); hence \(M(t)\) bounds the number of cells at edge semiperimeter \(t\).
\item[(b)] For every \(t\ge3\) there is a directionally convex polyhex \(K\) with
\(|K|=M(t)\) and \(\sigma(K)\le t\).
\item[(c)] The minimum edge perimeter of an \(N\)-cell polyhex is
\[
   p_{\min}(N)=2g(N),\qquad g(N)=\left\lceil\sqrt{12N-3}\,\right\rceil ,
\]

\item[(d)] Every \(N\)-cell minimizer \(S\) is directionally convex and satisfies
\(\sigma(S)=\lambda(S)=g(N)\).
\item[(e)] Writing \(\exc(S):=p(S)-p_{\min}(|S|)\), a polyhex \(S\) satisfies
\(\exc(S)\le2\) if and only if \(|S|>M(\sigma(S)-2)\).
\end{enumerate}
\end{theorem}

Parts (a)--(d) record classical or previously established edge-isoperimetric
quantities and properties, alongside self-contained proofs used in this note.
In particular, part (c) is the Harary--Harborth minimum edge-perimeter
formula \cite{HararyHarborth1976}. Yang's capacity results
\cite[Theorems~7--9 and Section~10]{Yang2002} include a
three-direction slice-count bound and the generalized capacity
\(\operatorname{round}(\lfloor p/2\rfloor^{2}/12)
=M(\lfloor p/2\rfloor)\) for a perimeter budget \(p\).
Consequently the capacity function \(M\) is not new here; the new
contribution is the controlled one-cell deletion proof below.

The directional convexity in (d) is a known result: Davoli, Piovano and Stefanelli
\cite[Proposition~3.3]{DavoliPiovanoStefanelli2017} prove that every minimizer of the
edge-isoperimetric problem on the triangular lattice is \(3\)-convex, which is
directional convexity in the present terminology; the proof below is the self-contained
one of the main paper, via the hull identities.

\begin{proof}
(a) Let \(H=H(S)\). By Lemma~\ref{lem:hull} and Lemma~\ref{lem:hex}(v),
\[
   12|S|\le12|H|=\lambda(H)^{2}+3-\bigl(3\delta(H)^{2}+2Q(H)\bigr)\le\lambda(S)^{2}+3
   \le\sigma(S)^{2}+3 ,
\]
the last step by Lemma~\ref{lem:runs}. Since \(|S|\) is an integer, \(|S|\le M(\sigma(S))\).

(b) For \(t\equiv3\pmod 6\), say \(t=6m+3\), take \(K=H(3m+1;m,m,m)\); for \(t=6m\) take
\(H(3m-1;m-1,m-1,m-1)\); for \(t=6m+1\) take \(H(3m;m,m,m-1)\); for \(t=6m+2\) take
\(H(3m+1;m,m,m+1)\); for \(t=6m+4\) take \(H(3m+1;m,m,m-1)\); for \(t=6m+5\) take
\(H(3m+2;m,m,m+1)\). In each case \eqref{eq:adm} holds for the indicated range of \(m\),
Lemma~\ref{lem:hex}(ii) gives \(\sigma(K)=t\), and Lemma~\ref{lem:hex}(iv),(v) give
\(12|K|=t^{2}+3-D\) with \(D\in\{0,3,4,7\}\) according to the residue of \(t\), so
\(|K|=M(t)\). The single case not covered by an admissible triple is \(t=4\); there
\(M(4)=1\) and a single cell has edge semiperimeter \(3\le4\).

(c) Put \(t_{N}:=\min\{t\ge3:\ M(t)\ge N\}\). By (a) every \(N\)-cell polyhex has
\(M(\sigma(S))\ge N\), hence \(\sigma(S)\ge t_{N}\). Conversely (b) produces a
directionally convex \(K\) with \(|K|=M(t_{N})\ge N\) and \(\sigma(K)\le t_{N}\), and
Corollary~\ref{cor:subsets} extracts from it an \(N\)-cell polyhex of edge semiperimeter at
most \(t_{N}\). Therefore \(p_{\min}(N)=2t_{N}\). Finally, for integers \(t\),
\[
   M(t)\ge N\iff \frac{t^{2}+3}{12}\ge N\iff t\ge\sqrt{12N-3},
\]
and \(\sqrt{12N-3}\ge3\) for \(N\ge1\), so \(t_{N}=g(N)\).

(d) By (c), \(\sigma(S)=g(N)\). Let \(H=H(S)\). Applying (a) to \(H\), which is
directionally convex with \(\sigma(H)=\lambda(H)\), gives \(M(\lambda(H))\ge|H|\ge N\) and
therefore \(\lambda(H)\ge t_{N}=g(N)\). Combining with Lemmas~\ref{lem:runs}
and~\ref{lem:hull},
\[
   g(N)=\sigma(S)=R(S)\ \ge\ \lambda(S)=\lambda(H)\ \ge\ g(N),
\]
so \(R(S)=\lambda(S)\) and \(S\) is directionally convex.

(e) Write \(t=\sigma(S)\) and \(k=|S|\); by (c), \(\exc(S)=2\bigl(t-g(k)\bigr)\). Since
both sides are integers, \(\exc(S)>2\) is equivalent to \(g(k)\le t-2\). If \(t-2\ge3\),
then by the definition of \(g\) as \(t_{k}\) and the monotonicity of \(M\) this is
equivalent to \(M(t-2)\ge k\). If \(t-2<3\), then \(t\in\{3,4\}\) because a polyhex has
\(t\ge3\); here \(M(t-2)=0<k\) and \(g(k)\ge3>t-2\), so both statements are false.
In both cases \(\exc(S)>2\) is equivalent to \(M(t-2)\ge k\), which is the assertion.
\end{proof}

\begin{remark}\label{rem:attain}
Parts (a) and (b) together say exactly that
\[
   M(t)=\max\{\,|K|:\ K\text{ a polyhex with }\sigma(K)\le t\,\},
\]
in words: \(M(t)\) is the largest number of cells that fit within an edge-semiperimeter
budget of \(t\). The budget form is the one the rest of the paper uses, and the
inequality in (b) is not a weakening: for every \(t\ge3\) other than \(t=4\) the
hexagons listed in the proof do satisfy \(\sigma(K)=t\), and the value \(t=4\) is not the
edge semiperimeter of any polyhex at all. Indeed a polyhex \(K\) with \(\sigma(K)=4\)
would have \(|K|\le M(4)=1\) by (a), while the one-cell polyhex has edge semiperimeter
\(3\).
\end{remark}

\subsection{The plateau bridge}\label{sec:2.7}

The function \(N\mapsto g(N)\) is nondecreasing and takes each of its values on a block of
consecutive integers. For \(\gamma\) in the range of \(g\) let
\[
   N_{b}(\gamma):=\min\{\,N\ge1:\ g(N)=\gamma\,\}
\]
be the first cell count at which the minimum edge semiperimeter equals \(\gamma\).

\begin{lemma}\label{lem:bridge}
For every \(\gamma\) in the range of \(g\),
\begin{equation}\label{eq:bridge}
   N_{b}(\gamma)=M(\gamma-1)+1 ,
\end{equation}
This identifies the first size in the plateau of threshold \(\gamma\).
\end{lemma}

\begin{proof}
For integers \(N\ge1\) and \(t\ge3\) we have, as in the proof of
Theorem~\ref{thm:capacity}(c), \(g(N)\le t\) if and only if \(N\le M(t)\). Hence
\(\{N:\ g(N)\le\gamma-1\}=\{1,\dots,M(\gamma-1)\}\), so the least \(N\) with
\(g(N)\ge\gamma\) is \(M(\gamma-1)+1\); as \(\gamma\) lies in the range of \(g\), that
least \(N\) satisfies \(g(N)=\gamma\).
\end{proof}




\section{Corridor and admissibility}\label{sec:corridor}

The main paper proves the growth theorem by deleting a whole extreme line at a time and
staying inside the minimizers. The proof recorded here descends one cell at a time and
controls the result through a pair of competing budgets. It reads its geometry from the
framework above, and it shares with the main paper the rigidity statement at exceptional
plateau starts (Proposition~\ref{pro:plateau} below, of which the main paper uses the two
exceptional cases).

\subsection{The lower bound on the line count}\label{sec:9.1}

\begin{lemma}\label{lem:lowerlines}
Every polyhex \(S\) satisfies \(\lambda(S)\ge g(|S|)\).
\end{lemma}

\begin{proof}
Let \(H=H(S)\). By Lemma~\ref{lem:hex}(ii) the hull is directionally convex with
\(\sigma(H)=\lambda(H)\), so Theorem~\ref{thm:capacity}(a) applied to \(H\) gives
\(|H|\le M(\lambda(H))\). By Lemma~\ref{lem:hull}, \(S\subseteq H\) and
\(\lambda(H)=\lambda(S)\); hence \(|S|\le M(\lambda(S))\), and the definition of \(g\) as
the least \(t\) with \(M(t)\ge|S|\) gives \(\lambda(S)\ge g(|S|)\).
\end{proof}

In words: the number of occupied lines can never fall below the threshold, because the hull
would then be too small to hold the configuration.

\subsection{Excess, band and corridor}\label{sec:9.2}

\begin{definition}\label{def:eps}
For a polyhex \(S\) the \emph{line-count excess} is
\[
   \varepsilon(S):=\lambda(S)-g(|S|)\ \ge 0 ,
\]
in words: the number of occupied lines above the threshold. We set
\[
   \mathcal{B}_{2}:=\{\,S:\exc(S)\le2\,\},\qquad
   \mathcal{C}:=\{\,S:\ S\ \text{directionally convex},\ \varepsilon(S)\in\{0,1\}\,\},
\]
and call \(\mathcal{B}_{2}\) the \emph{2-excess band} and \(\mathcal{C}\) the \emph{unit
line-excess corridor}, or simply the \emph{corridor}.
\end{definition}

If \(S\) is directionally convex then \(p(S)=2\lambda(S)\) by Lemma~\ref{lem:runs}, and
\(p_{\min}=2g\) by Theorem~\ref{thm:capacity}(c), so
\begin{equation}\label{eq:exc2eps}
   \exc(S)=2\,\varepsilon(S),
\end{equation}
in words: for a directionally convex configuration the edge-perimeter excess is exactly
twice the line-count excess. Hence \(\mathcal{C}\subseteq\mathcal{B}_{2}\), and the
inclusion is strict: the band also contains configurations that are not directionally
convex, an explicit nine-cell one being given in Section~\ref{sec:counter}.

\subsection{Thresholds, plateaux and budgets}\label{sec:9.3}

The threshold \(g\) is nondecreasing; write its increment as the \emph{threshold jump}
\[
   j(N):=g(N)-g(N-1).
\]
For \(N\ge3\) one has \(j(N)\in\{0,1\}\). (The value \(j(2)=5-3=2\) is exceptional; the
iteration below is carried out for sizes at least \(4\).)

The function \(g\) is not surjective: \(g(1)=3\) and \(g(2)=5\), so \(\gamma=4\) is not a
value of \(g\). For \(\gamma\) in the range of \(g\) we use the plateau start
\(N_{b}(\gamma)=\min\{n:g(n)=\gamma\}\) of Section~\ref{sec:2.7}, for which
Lemma~\ref{lem:bridge} gives the closed form
\begin{equation}\label{eq:nbform}
   N_{b}(\gamma)=M(\gamma-1)+1 .
\end{equation}
How far \(S\) has advanced along its plateau is the \emph{plateau depth}
\[
   \tau(S):=|S|-N_{b}\bigl(g(|S|)\bigr)\ \ge0 .
\]
(The symbol \(\tau\) is used in this part for the plateau depth only; edge semiperimeters
are written \(\sigma\) throughout, as in the framework above.)

For a directionally convex polyhex we write \(r_{1},\dots,r_{6}\) for the support traces
and \(\rho(S)=\min_{i}r_{i}\) as in Definition~\ref{def:traces}; \(\rho(S)\ge1\) because no
extreme line is empty. Every configuration to which \(\rho\) is applied below is verified
to be directionally convex first, as Definition~\ref{def:traces} requires.

\subsection{The effect of a single deletion}\label{sec:9.4}

Let \(c\in S\) and \(S':=S\sm\{c\}\). Let \(\solo(S,c)\) denote the \emph{line-loss count},
the number of occupied lines that disappear when \(c\) is deleted. No new line is occupied,
so \(\lambda(S')=\lambda(S)-\solo(S,c)\), and Definition~\ref{def:eps} gives at once
\begin{equation}\label{eq:delta1}
   \varepsilon(S')=\varepsilon(S)+j(|S|)-\solo(S,c),
\end{equation}
in words: deleting a cell raises the line-count excess by the threshold jump and lowers it
by the number of lines vacated. This identity uses neither connectedness nor convexity. If
\(S\) is connected with \(|S|\ge2\) then \(c\) has a neighbour, and each of the six
directions of adjacency preserves one of the three keys, so \(\solo(S,c)\le2\).

\subsection{The admissibility condition}\label{sec:adm}

We wish to keep \(\varepsilon\le1\) at every step. By \eqref{eq:delta1} with
\(\solo\ge0\) and \(j\le1\), the value \(\varepsilon(S')\) can reach \(2\) only when
\[
   \varepsilon(S)=1,\qquad j(|S|)=1,\qquad \solo(S,c)=0
\]
all hold at once: one spare line is already in use, the threshold drops as the size
decreases, and no line is vacated to pay for it. The invariant of this part exists only to
forbid that combination.

The device is a comparison of two budgets. The plateau depth \(\tau(S)\) says how many
further deletions are possible before the threshold drops again. The shortest support trace
\(\rho(S)\) says how many deletions are needed to exhaust the shortest extreme line, one
cell at a time; at the deletion that exhausts it the line itself disappears, so
\(\solo\ge1\) and the term \(+j\) in \eqref{eq:delta1} is cancelled. So
\(\rho(S)\le\tau(S)\) is exactly the statement that \emph{the line runs out before the
threshold drops}.

\begin{definition}[admissibility condition]\label{def:adm}
For a directionally convex polyhex \(S\),
\[
   \Adm(S)\quad:\Longleftrightarrow\quad
   \varepsilon(S)=0\quad\text{or}\quad\bigl(\varepsilon(S)=1\ \text{and}\ \rho(S)\le\tau(S)\bigr).
\]
\end{definition}

\begin{lemma}\label{lem:admjump}
If \(\Adm(S)\) then \(\varepsilon(S)=1\) and \(j(|S|)\ge1\) cannot hold simultaneously.
\end{lemma}

\begin{proof}
Assume \(j(|S|)\ge1\). As \(g\) is nondecreasing, \(n<|S|\) gives
\(g(n)\le g(|S|-1)<g(|S|)\), so \(|S|\) is itself the least size at which \(g\) takes the
value \(g(|S|)\), that is \(\tau(S)=0\). If \(\varepsilon(S)=1\), the second alternative of
\(\Adm\) would give \(\rho(S)\le0\), contradicting \(\rho(S)\ge1\).
\end{proof}

\begin{lemma}[initialization]\label{lem:init}
Every edge-perimeter minimizer \(S\) satisfies \(S\in\mathcal{C}\) and \(\Adm(S)\).
\end{lemma}

\begin{proof}
By Theorem~\ref{thm:capacity}(d), \(S\) is directionally convex with
\(\lambda(S)=g(|S|)\), that is \(\varepsilon(S)=0\). The first alternative of
Definition~\ref{def:adm} applies.
\end{proof}

What remains is to show that such a deletion exists at all
(Section~\ref{sec:legal}) and that it preserves \(\Adm\) (Section~\ref{sec:preserve}).
These two are established separately and are not combined into a single statement; the
separation is the point of this proof, and Section~\ref{sec:counter} exhibits a configuration for which
existence holds while admissibility fails.

\section{Legal one-cell deletion}\label{sec:legal}

\begin{proposition}\label{pro:legal}
Let \(S\in\mathcal{C}\) with \(|S|\ge4\). Then there is a cell \(c\) which is an endpoint of
a shortest extreme line of \(S\) and whose deletion leaves \(S\) connected.
\end{proposition}

Such a cell is called a \emph{legal deletion}. Note that a single cell is removed; this is
not the block peeling of the main paper, and the two operations are not interchangeable.

\subsection{The witness cell}\label{sec:11.1}

Among the six extreme lines of \(S\) choose one attaining \(\rho(S)\); if there are several,
any one will do. The symmetry group of the lattice permutes the three families cyclically
and exchanges the minimum and maximum side within a family, so we may assume the chosen line
is \(q=q_{\max}\). Let
\[
   c^{*}:=\operatorname*{argmax}\{\,r:(q_{\max},r)\in S\,\}
\]
be the cell of that line with the largest \(r\); we call it the \emph{witness cell}. It is
an endpoint of the trace on the chosen extreme line.

\subsection{Local structure at the witness cell}\label{sec:11.2}

\begin{lemma}\label{lem:local}
None of \(c^{*}+d_{0}\), \(c^{*}+d_{1}\), \(c^{*}+d_{5}\) belongs to \(S\). Consequently
\(\deg_{S}(c^{*})\le3\), and the neighbours of \(c^{*}\) inside \(S\) are confined to the
three directions \(d_{2},d_{3},d_{4}\), which are consecutive in the cyclic order.
\end{lemma}

\begin{proof}
Both \(d_{0}=(1,0)\) and \(d_{5}=(1,-1)\) increase \(q\), and \(c^{*}\) lies on
\(q=q_{\max}\). The direction \(d_{1}=(0,1)\) preserves \(q\) and increases \(r\), and
\(c^{*}\) has the largest \(r\) on that line.
\end{proof}

\begin{corollary}\label{cor:endpoint3}
In each of the three families, \(c^{*}\) is an endpoint of the trace on the line containing
it.
\end{corollary}

\begin{proof}
In the \(q\)-family the line is \(q=q_{\max}\) and \(c^{*}\) attains the maximum of \(r\) on
it. In the \(r\)- and \(s\)-families the free coordinate along a line is \(q\), and
\(c^{*}\) has \(q=q_{\max}\), the largest value of \(q\) on all of \(S\).
\end{proof}

\begin{corollary}\label{cor:art}
If \(c^{*}\) is an articulation cell of \(S\), then \(\deg_{S}(c^{*})=2\) and its two
neighbours lie in the directions \(d_{2}\) and \(d_{4}\).
\end{corollary}

\begin{proof}
If \(\deg_{S}(c^{*})\le1\) then \(S\sm\{c^{*}\}\) is connected. If \(\deg_{S}(c^{*})=3\) the
neighbours are \(c^{*}+d_{2},c^{*}+d_{3},c^{*}+d_{4}\), and since \(d_{2},d_{3}\) and
\(d_{3},d_{4}\) are adjacent pairs of directions, these three cells are joined to one
another inside \(S\sm\{c^{*}\}\), so any path through \(c^{*}\) can be rerouted. If
\(\deg_{S}(c^{*})=2\) with directions \(\{d_{2},d_{3}\}\) or \(\{d_{3},d_{4}\}\) the same
applies. Only \(\{d_{2},d_{4}\}\) remains.
\end{proof}

\subsection{A degree-two articulation cell forces excess at least four}\label{sec:11.3}

\begin{lemma}[arithmetic lemma]\label{lem:arith}
Let \(a,b\ge1\) be integers with \((a,b)\neq(1,1)\). Then
\(g(a)+g(b)\ge g(a+b+1)+1\). For \((a,b)=(1,1)\) the statement fails.
\end{lemma}

\begin{proof}
Put \(A:=g(a)\), \(B:=g(b)\). By definition \(A\ge\sqrt{12a-3}\), that is
\(12a\le A^{2}+3\), and likewise for \(b\); hence
\(12(a+b+1)-3\le A^{2}+B^{2}+15\), and by monotonicity of the ceiling of the square root,
\(g(a+b+1)\le\lceil\sqrt{A^{2}+B^{2}+15}\,\rceil\). Since \(A+B-1\) is a positive integer,
\(\lceil\sqrt{A^{2}+B^{2}+15}\,\rceil\le A+B-1\) is equivalent to
\(A^{2}+B^{2}+15\le(A+B-1)^{2}\), that is to \((A-1)(B-1)\ge8\). Now \(g\) is nondecreasing
with \(g(1)=3\) and \(g(2)=5\), so \(A=3\) if \(a=1\) and \(A\ge5\) if \(a\ge2\), and
similarly for \(B\); unless \(a=b=1\) we get \((A-1)(B-1)\ge2\cdot4=8\). Finally
\(A=B=3\) happens exactly for \(a=b=1\), where \(g(1)+g(1)=6=g(3)\).
\end{proof}

\begin{lemma}[Bridge-Excess Lemma]\label{lem:bridge-exc}
Let \(X\) be a polyhex with \(|X|=N\ge4\) having a degree-two articulation cell. Then
\(\exc(X)\ge4\).
\end{lemma}

\begin{proof}
Let \(c\) be such a cell. Every component of \(X\sm\{c\}\) contains a neighbour of \(c\),
and \(c\) has two neighbours, so there are at most two components; being an articulation
cell there are at least two. Call them \(X_{1},X_{2}\) and put \(n_{i}=|X_{i}|\), so
\(n_{1}+n_{2}=N-1\). Counting adjacent pairs, \(a(X)=a(X_{1})+a(X_{2})+2\), whence
\[
   p(X)=6N-2a(X)=p(X_{1})+p(X_{2})+2 .
\]
In words: the two pieces keep their own perimeters and the bridge contributes two further
boundary edges. With \(p(X_{i})\ge2g(n_{i})\) and Lemma~\ref{lem:arith},
\[
   p(X)\ \ge\ 2g(n_{1})+2g(n_{2})+2\ \ge\ 2\bigl(g(N)+1\bigr)+2\ =\ p_{\min}(N)+4 .
\]
Lemma~\ref{lem:arith} needs \((n_{1},n_{2})\neq(1,1)\), which would mean \(N=3\).
\end{proof}

\subsection{Proof of Proposition~\ref{pro:legal}}\label{sec:11.4}

\begin{proof}
Let \(S\in\mathcal{C}\) with \(|S|\ge4\). By \eqref{eq:exc2eps}, \(\exc(S)\le2\). Choose the
witness cell \(c^{*}\). If it were an articulation cell then by Corollary~\ref{cor:art} it
would be one of degree two, and Lemma~\ref{lem:bridge-exc} would give \(\exc(S)\ge4\). Hence
\(S\sm\{c^{*}\}\) is connected, and \(c^{*}\) is a legal deletion.
\end{proof}

\begin{remark}[Role of the shortest-line hypothesis]\label{rem:shortest}
The argument just given used the shortest extreme line only to choose \(c^{*}\); the
connectedness conclusion did not depend on that choice. Indeed a theorem
(recorded in a separate note) shows, by the same local argument, that for \(|S|\ge4\)
and \(\exc(S)\le2\) \emph{no} endpoint of \emph{any} extreme line is an articulation cell,
whether or not \(S\) is directionally convex. The restriction to a shortest extreme line is needed for a
different reason: it is what the budget bookkeeping of Section~\ref{sec:preserve} tracks
through \(\rho(S)\). The two hypotheses of that theorem cannot be dropped: in
the bent three-cell polyhex \(\{(0,1),(0,2),(1,0)\}\), of excess \(2\), the cell \((0,1)\)
is an endpoint of the extreme line \(q=0\) and its deletion disconnects the set; and in the
four-cell polyhex \(\{(0,1),(0,2),(0,3),(1,0)\}\), of excess \(4\), the same is true of
\((0,1)\).
\end{remark}

\section{The plateau-start bound}\label{sec:plateau}

\begin{proposition}[Plateau-Start Bound]\label{pro:plateau}
Let \(\gamma\) be in the range of \(g\) with \(\gamma\ge7\), and let \(S\) be an
edge-perimeter minimizer with \(|S|=N_{b}(\gamma)\). Then
\[
   \rho(S)\ \le\ N_{b}(\gamma)-N_{b}(\gamma-1).
\]
\end{proposition}

Informally: at the start of a plateau the shortest extreme line of a minimizer cannot be
longer than the distance down to the start of the previous plateau. This is what carries
admissibility across the moment at which the threshold drops.

\subsection{Two special values}\label{sec:12.1}

By \eqref{eq:nbform}, for an integer \(h\ge2\),
\begin{equation}\label{eq:nbspecial}
   N_{b}(6h-3)=3h^{2}-4h+2,\qquad N_{b}(6h-1)=3h^{2}-2h+1 ,
\end{equation}
since \(\bigl((6h-4)^{2}+3\bigr)/12\) has integer part \(3h^{2}-4h+1\) and
\(\bigl((6h-2)^{2}+3\bigr)/12\) has integer part \(3h^{2}-2h\).

\subsection{Excess of the hull sides}\label{sec:12.2}

All hull material is taken from the framework above: the parametrization \(H(L;u,v,w)\) of
Definition~\ref{def:hexagon}, the side lines \(\Lambda_{1},\dots,\Lambda_{6}\) in cyclic
order, the side lengths \(\ell_{1},\dots,\ell_{6}\) and corners \(\kappa_{i}\) of
Lemma~\ref{lem:hex}(i), the sum rule \eqref{eq:sidesum}, the opposite-side difference
\(\delta(H)\) of \eqref{eq:delta}, and the shortest side length \(h=\min_{i}\ell_{i}\).
Write
\[
   x_{i}:=\ell_{i}-h\ \ge0
\]
for the \emph{side excess}, the amount by which a hull side exceeds the shortest hull side.


\begin{lemma}[Opposite-Pair Excess Lemma]\label{lem:oppair}
If \(\sum_{i}x_{i}\le2\) then \(x_{i}=x_{i+3}\) for every \(i\); in particular
\(\sum_{i}x_{i}\) is even, so \(\sum_{i}x_{i}=1\) is impossible, and if
\(\sum_{i}x_{i}=2\) then exactly one pair of opposite sides has length \(h+1\) and the other
four sides have length \(h\).
\end{lemma}

\begin{proof}
Since the \(x_{i}\) differ from the \(\ell_{i}\) by the constant \(h\), \eqref{eq:delta}
gives \(x_{1}-x_{4}=x_{3}-x_{6}=x_{5}-x_{2}=\delta(H)\). If \(\delta(H)\ge1\) then
\(x_{1},x_{3},x_{5}\ge1\) and \(\sum x_{i}\ge3\); if \(\delta(H)\le-1\) the same holds for
\(x_{2},x_{4},x_{6}\). Both contradict the hypothesis, so \(\delta(H)=0\) and
\(\sum x_{i}=2(x_{1}+x_{3}+x_{5})\).
\end{proof}

\begin{lemma}[a complete boundary forces equality with the hull]\label{lem:complete}
Let \(S\) be a directionally convex polyhex and \(H=H(S)\). If \(H\cap\Lambda_{i}\subseteq S\)
for all six \(i\), then \(S=H\).
\end{lemma}

\begin{proof}
Fix a value \(q_{0}\). By Definition~\ref{def:hexagon} the cells of \(H\) on the line
\(q=q_{0}\) are those with \(r\) in the interval
\(\bigl[\max(0,w-q_{0}),\ \min(B-1,L-1-q_{0})\bigr]\). Its lower end has \(r=0\) or
\(q+r=w\), so it lies on \(\Lambda_{5}\) or \(\Lambda_{6}\); its upper end has \(r=B-1\) or
\(q+r=L-1\), so it lies on \(\Lambda_{2}\) or \(\Lambda_{3}\). By hypothesis both ends lie
in \(S\), and the trace of \(S\) on this line is a run, so it contains everything between
them. As \(q_{0}\) was arbitrary, \(H\subseteq S\).
\end{proof}

\subsection{Proof of Proposition~\ref{pro:plateau}}\label{sec:12.3}

\begin{proof}
Let \(N=N_{b}(\gamma)\) with \(\gamma\ge7\). By Theorem~\ref{thm:capacity}(d), \(S\) is
directionally convex with \(\lambda(S)=\gamma\), and by Lemma~\ref{lem:hull}(ii)
\(\lambda(H)=\gamma\) for \(H=H(S)\). Since \(r_{i}\le\ell_{i}\) for each \(i\), we have
\(\rho(S)\le h\). The minimum of six numbers is at most their mean, and it is an integer, so
\eqref{eq:sidesum} gives
\begin{equation}\label{eq:betadef}
   \rho(S)\ \le\ h\ \le\ \Bigl\lfloor\frac{\sum_{i}\ell_{i}}{6}\Bigr\rfloor
   =\Bigl\lfloor\frac{\lambda(H)+3}{6}\Bigr\rfloor
   =\Bigl\lfloor\frac{\gamma+3}{6}\Bigr\rfloor=:\beta(\gamma).
\end{equation}

First, \(N_{b}(\gamma-1)\) is defined: \(N\ge4\) and \(N\) is the least size with
\(g=\gamma\), so \(g(N-1)<\gamma\); since \(j(N)\in\{0,1\}\) for \(N\ge3\) we get
\(g(N-1)=\gamma-1\), so \(\gamma-1\) is in the range of \(g\). Write \(\gamma-1=6k+t\) with
\(t\in\{0,\dots,5\}\). By \eqref{eq:nbform},
\(N_{b}(\gamma)-1=M(\gamma-1)=3k^{2}+kt+\lfloor(t^{2}+3)/12\rfloor\), and
\(\lfloor(t^{2}+3)/12\rfloor\) equals \(0,0,0,1,1,2\) for \(t=0,\dots,5\). Taking
differences, and noting that \(\gamma-2=6k+(t-1)\) when \(t\ge1\) while
\(\gamma-2=6(k-1)+5\) when \(t=0\), gives the following table.

\begin{center}
\begin{tabular}{rrrrr}
\hline
\(t\) & \(\gamma\bmod 6\) & \(N_{b}(\gamma)-N_{b}(\gamma-1)-k\) & \(\beta(\gamma)-k\)
 & \(\beta(\gamma)-\bigl(N_{b}(\gamma)-N_{b}(\gamma-1)\bigr)\) \\
\hline
0 & 1 & 0 & 0 & 0 \\
1 & 2 & 0 & 0 & 0 \\
2 & 3 & 0 & 1 & \textbf{1} \\
3 & 4 & 1 & 1 & 0 \\
4 & 5 & 0 & 1 & \textbf{1} \\
5 & 0 & 1 & 1 & 0 \\
\hline
\end{tabular}
\end{center}

Hence \(\beta(\gamma)-\bigl(N_{b}(\gamma)-N_{b}(\gamma-1)\bigr)=1\) exactly when
\(\gamma\equiv3\) or \(5\pmod 6\), and \(\beta(\gamma)=N_{b}(\gamma)-N_{b}(\gamma-1)\)
otherwise.

\medskip\noindent\textbf{Case 1: \(\gamma\not\equiv3,5\pmod6\).} Then
\(\beta(\gamma)=N_{b}(\gamma)-N_{b}(\gamma-1)\) and \eqref{eq:betadef} gives the claim at
once.

\medskip\noindent\textbf{Case 2: \(\gamma\equiv3\pmod6\).} Write \(\gamma=6h-3\); from
\(\gamma\ge7\) we get \(h\ge2\). Here \(\beta(\gamma)=h\) and we must show
\(\rho(S)\le h-1\). Suppose \(\rho(S)=h\). Then \eqref{eq:betadef} forces
\(h=\min_{i}\ell_{i}=\beta(\gamma)\), while \(\sum_{i}\ell_{i}=\gamma+3=6h\); as all six
sides are at least \(h\), all six equal \(h\). Solving
\(\ell_{2}=v+1=\ell_{4}=u+1=\ell_{6}=w+1=h\) and \(\ell_{1}=L-v-w=h\) gives
\[
   u=v=w=h-1,\qquad L=3h-2,
\]
so by Lemma~\ref{lem:hex}(iv)
\[
   |H|=T(3h-2)-3\,T(h-1)=3h^{2}-3h+1 .
\]
By \eqref{eq:nbspecial}, \(N=3h^{2}-4h+2\), so \(|H|-N=h-1\ge1\). But \(\rho(S)=h\) forces
\(h\le r_{i}\le\ell_{i}=h\) for every \(i\), that is, every side of the hull lies entirely
in \(S\); Lemma~\ref{lem:complete} then gives \(S=H\), contradicting \(|S|=N<|H|\).

\medskip\noindent\textbf{Case 3: \(\gamma\equiv5\pmod6\).} Write \(\gamma=6h-1\), again
\(h\ge2\). Here \(\beta(\gamma)=h\) and we must show \(\rho(S)\le h-1\). Suppose
\(\rho(S)=h\); as before \(\min_{i}\ell_{i}=h\), while now
\(\sum_{i}\ell_{i}=\gamma+3=6h+2\), so \(\sum_{i}x_{i}=2\). By Lemma~\ref{lem:oppair}
exactly one pair of opposite sides has length \(h+1\). The three families are permuted
cyclically by the \(120^{\circ}\) rotation of the lattice, so we may assume that pair is
\((\ell_{1},\ell_{4})\), the two sides of the \(q\)-family. Solving
\(\ell_{4}=u+1=h+1\), \(\ell_{2}=v+1=h\), \(\ell_{6}=w+1=h\) and \(\ell_{1}=L-v-w=h+1\)
gives
\[
   u=h,\qquad v=w=h-1,\qquad L=3h-1
\]
(one checks \(\ell_{3}=\ell_{5}=h\) and \(\lambda(H)=3L-(u+v+w)=6h-1=\gamma\)), and
\[
   |H|=T(3h-1)-T(h)-2\,T(h-1)=3h^{2}-h .
\]
By \eqref{eq:nbspecial}, \(N=3h^{2}-2h+1\), so again \(|H|-N=h-1\ge1\).

From \(\rho(S)=h\) it follows that the four sides of length \(h\), namely
\(\ell_{2},\ell_{3},\ell_{5},\ell_{6}\), lie entirely in \(S\). The long side
\(\Lambda_{1}\) is adjacent in the cyclic order to \(\Lambda_{6}\) and \(\Lambda_{2}\),
with which it shares the corners \(\kappa_{6}\) and \(\kappa_{1}\); both corners therefore
lie in \(S\), and since the trace of \(S\) on \(\Lambda_{1}\) is a run containing both ends,
\(\Lambda_{1}\) lies entirely in \(S\) as well. The opposite side \(\Lambda_{4}\) is handled
the same way through its neighbours \(\Lambda_{3}\) and \(\Lambda_{5}\). So all six sides
lie in \(S\), and Lemma~\ref{lem:complete} gives \(S=H\), contradicting \(|S|=N<|H|\).
\end{proof}

\subsection{The form used in Section~\ref{sec:preserve}}\label{sec:12.4}

\begin{corollary}\label{cor:plateau}
Let \(S\in\mathcal{C}\) with \(|S|\ge4\), \(\varepsilon(S)=0\) and \(j(|S|)=1\). Put
\(\gamma:=g(|S|)\). Then \(|S|=N_{b}(\gamma)\) and
\(\rho(S)\le N_{b}(\gamma)-N_{b}(\gamma-1)\).
\end{corollary}

\begin{proof}
As \(S\) is directionally convex with \(\varepsilon(S)=0\), Lemma~\ref{lem:runs} and
Theorem~\ref{thm:capacity}(c) give \(p(S)=2\lambda(S)=2g(|S|)=p_{\min}(|S|)\), so \(S\) is a
minimizer. From \(j(|S|)=1\) and monotonicity of \(g\) we get \(g(n)<\gamma\) for
\(n<|S|\), that is \(|S|=N_{b}(\gamma)\); and \(|S|\ge4\) gives \(\gamma\ge g(4)=7\). Apply
Proposition~\ref{pro:plateau}.
\end{proof}

\section{Preservation and the alternative proof}\label{sec:preserve}

\begin{proposition}\label{pro:preserve}
Let \(S\in\mathcal{C}\) with \(|S|\ge4\) and \(\Adm(S)\), let \(c^{*}\) be the witness cell
of Section~\ref{sec:11.1}, and put \(S':=S\sm\{c^{*}\}\). Then \(S'\in\mathcal{C}\) and
\(\Adm(S')\).
\end{proposition}

Connectedness of \(S'\) is Proposition~\ref{pro:legal}. What remains is directional
convexity, \(\varepsilon(S')\in\{0,1\}\), and \(\Adm(S')\).

\subsection{Directional convexity is preserved}\label{sec:13.1}

\begin{lemma}\label{lem:convpres}
Let \(S\) be directionally convex and let \(c\in S\) be, in each of the three families, an
endpoint of the trace on the line containing it. Then \(S\sm\{c\}\) is directionally convex.
\end{lemma}

\begin{proof}
In each family exactly one line has its trace changed. An unchanged trace is still a run. The
changed trace is a run of which \(c\) is an end, so removing \(c\) leaves a run, possibly
empty.
\end{proof}

By Corollary~\ref{cor:endpoint3} the witness cell is an endpoint in all three families, so
\(S'\) is directionally convex. In particular \(\rho(S')\) is defined, which is what permits
its use below.

\subsection{Tracking the two budgets}\label{sec:13.2}

Write \(\gamma:=g(|S|)\).

\begin{lemma}[plateau depth after a deletion]\label{lem:taudiff}
\textup{(i)} If \(j(|S|)=0\) then \(\tau(S')=\tau(S)-1\).
\textup{(ii)} If \(j(|S|)=1\) then \(|S|=N_{b}(\gamma)\), \(\tau(S)=0\) and
\(\tau(S')=N_{b}(\gamma)-1-N_{b}(\gamma-1)\).
\end{lemma}

\begin{proof}
(i) Here \(g(|S|-1)=\gamma\), so \(\tau(S')=(|S|-1)-N_{b}(\gamma)=\tau(S)-1\), and
\(N_{b}(\gamma)\le|S|-1\) gives \(\tau(S')\ge0\).
(ii) Monotonicity gives \(g(n)<\gamma\) for \(n<|S|\), so \(|S|=N_{b}(\gamma)\) and
\(\tau(S)=0\); also \(g(|S|-1)=\gamma-1\), whence
\(\tau(S')=(|S|-1)-N_{b}(\gamma-1)\).
\end{proof}

\begin{lemma}[shortest extreme line decrement]\label{lem:decr}
Let \(S\) be directionally convex, let \(c\in S\) lie on a shortest extreme line, and suppose
\(\rho(S)\ge2\) and that \(S\sm\{c\}\) is directionally convex. Then
\(\rho(S\sm\{c\})=\rho(S)-1\).
\end{lemma}

\begin{proof}
First, the six extreme lines of \(S\) and of \(S'\) coincide as pairs (family, key). Fix a
family and let \(K\) be the set of keys occurring on \(S\). If the key of \(c\) is neither
the minimum nor the maximum of \(K\), removing \(c\) changes neither. If it is extremal, the
corresponding line is an extreme line and carries at least \(\rho(S)\ge2\) cells, so another
cell has the same key and the extremes are again unchanged. Next, for each extreme line
\(\Lambda\) we have \(|S'\cap\Lambda|=|S\cap\Lambda|-[\,c\in\Lambda\,]\). All six values are
at least \(\rho(S)\); the value on the shortest extreme line containing \(c\) drops to
\(\rho(S)-1\), and every other drops by at most one, so stays at least \(\rho(S)-1\).
\end{proof}

\begin{lemma}\label{lem:solo0}
If \(c\) lies on a shortest extreme line and \(\solo(S,c)=0\), then \(\rho(S)\ge2\).
\end{lemma}

\begin{proof}
If \(\rho(S)=1\) the shortest extreme line containing \(c\) consists of \(c\) alone, so
deleting \(c\) vacates that line and \(\solo(S,c)\ge1\).
\end{proof}

\subsection{The case distinction}\label{sec:13.3}

By \(\Adm(S)\) and Lemma~\ref{lem:admjump}, the combination \(\varepsilon(S)=1\) and
\(j(|S|)=1\) does not occur. Moreover \eqref{eq:delta1} together with
\(\varepsilon(S')\ge0\), which is Lemma~\ref{lem:lowerlines} applied to \(S'\), gives
\begin{equation}\label{eq:solobound}
   \solo(S,c^{*})\ \le\ \varepsilon(S)+j(|S|).
\end{equation}
With \(\varepsilon(S)\in\{0,1\}\) and \(j(|S|)\in\{0,1\}\) there are three cases.

\medskip\noindent\textbf{Case A: \(\varepsilon(S)=0\), \(j(|S|)=0\).} By
\eqref{eq:solobound}, \(\solo=0\), so \eqref{eq:delta1} gives \(\varepsilon(S')=0\) and
\(\Adm(S')\) holds by its first alternative.

\medskip\noindent\textbf{Case B: \(\varepsilon(S)=0\), \(j(|S|)=1\).} Here \(\solo\le1\).

\emph{B1: \(\solo=1\).} Then \(\varepsilon(S')=0+1-1=0\); first alternative.

\emph{B2: \(\solo=0\).} Then \(\varepsilon(S')=1\) and the second alternative must be
established. All hypotheses of Corollary~\ref{cor:plateau} hold, so
\(\rho(S)\le N_{b}(\gamma)-N_{b}(\gamma-1)\). By Lemma~\ref{lem:solo0},
\(\rho(S)\ge2\), and \(S'\) is directionally convex by Section~\ref{sec:13.1}, so
Lemma~\ref{lem:decr} gives \(\rho(S')=\rho(S)-1\). By Lemma~\ref{lem:taudiff}(ii),
\(\tau(S')=N_{b}(\gamma)-1-N_{b}(\gamma-1)\). Therefore
\[
   \rho(S')=\rho(S)-1\ \le\ \bigl(N_{b}(\gamma)-N_{b}(\gamma-1)\bigr)-1\ =\ \tau(S') .
\]

Note that the bound of Corollary~\ref{cor:plateau}, decreased by one, coincides
\emph{identically} with the plateau depth after the deletion. In the language of
Section~\ref{sec:adm}: what remains of the shortest-line budget allowed at the start of the
plateau, after one unit has been spent, is exactly the plateau budget afterwards. The form
of \(\Adm\) is chosen to make this branch close.

\medskip\noindent\textbf{Case C: \(\varepsilon(S)=1\), \(j(|S|)=0\).} Here \(\solo\le1\).

\emph{C1: \(\solo=1\).} Then \(\varepsilon(S')=1+0-1=0\); first alternative.

\emph{C2: \(\solo=0\).} Then \(\varepsilon(S')=1\). The second alternative of \(\Adm(S)\)
gives \(\rho(S)\le\tau(S)\); by Lemmas~\ref{lem:solo0} and \ref{lem:decr},
\(\rho(S')=\rho(S)-1\), and by Lemma~\ref{lem:taudiff}(i), \(\tau(S')=\tau(S)-1\). Hence
\(\rho(S')\le\tau(S')\). This case does not use Section~\ref{sec:plateau}.

\medskip
In every case \(\varepsilon(S')\in\{0,1\}\) and \(\Adm(S')\); with
Section~\ref{sec:13.1} and Proposition~\ref{pro:legal} this gives \(S'\in\mathcal{C}\).

\begin{remark}\label{rem:sizefour}
What Proposition~\ref{pro:preserve} constrains is the size \emph{before} the deletion, so it
applies at \(|S|=4\) and its conclusion holds for \(|S'|=3\); it cannot be applied again
with \(S'\) in the role of \(S\). This is where the iteration stops, and sizes at most
\(3\) are handled directly below.
\end{remark}

\subsection{The one-cell deletion proof of the growth theorem}\label{sec:altmain}

\begin{lemma}\label{lem:small}
Every polyhex \(X\) with \(|X|\le3\) satisfies \(\exc(X)\le2\).
\end{lemma}

\begin{proof}
For \(|X|=1\), \(p=6=2g(1)\). For \(|X|=2\) there is one adjacent pair, so
\(p=12-2=10=2g(2)\). For \(|X|=3\) the number of adjacent pairs is \(2\) or \(3\), so
\(p\in\{12,14\}\), and \(p_{\min}(3)=2g(3)=12\).
\end{proof}

\begin{lemma}\label{lem:leaf}
Every polyhex \(X\) with \(|X|\ge2\) has a cell whose deletion leaves \(X\) connected.
\end{lemma}

\begin{proof}
The adjacency graph of \(X\) is connected with at least two vertices; a leaf of any spanning
tree is not a cut vertex.
\end{proof}

\begin{theorem}[one-cell growth]\label{thm:altgrowth}
Let \(S\) be an \(N\)-cell edge-perimeter minimizer. Then there are polyhexes
\(S_{1}\subset\cdots\subset S_{N}=S\), \(|S_{k}|=k\), each obtained from the previous one by
adding an adjacent cell, with \(\exc(S_{k})\le2\) for every \(k\).
\end{theorem}
\begin{proof}
For \(N\le3\) apply Lemma~\ref{lem:leaf} repeatedly, removing one cell at a time while
keeping connectedness; every configuration on the way is a polyhex of size at most \(3\), so
(b) holds by Lemma~\ref{lem:small}, and (a) follows as in the reversal step below. Assume
\(N\ge4\).

\emph{Initialization.} By Lemma~\ref{lem:init}, \(S\in\mathcal{C}\) and \(\Adm(S)\). Put
\(T_{N}:=S\).

\emph{Descent.} Suppose \(T_{k}\in\mathcal{C}\), \(|T_{k}|=k\ge4\) and \(\Adm(T_{k})\). By
Proposition~\ref{pro:legal} there is an endpoint \(c_{k}\) of a shortest extreme line of
\(T_{k}\) with \(T_{k}\sm\{c_{k}\}\) connected; put \(T_{k-1}:=T_{k}\sm\{c_{k}\}\). By
Proposition~\ref{pro:preserve}, \(T_{k-1}\in\mathcal{C}\) and \(\Adm(T_{k-1})\). Repeating
down to \(k=4\) produces \(T_{3}\) with \(|T_{3}|=3\); by Remark~\ref{rem:sizefour} the
step at \(k=4\) is legitimate and the iteration stops there. Sizes at most \(3\) are
continued with Lemma~\ref{lem:leaf}: choose \(T_{3}\supset T_{2}\supset T_{1}\) preserving
connectedness.

\emph{Staying in the band.} For \(k\ge4\) we have \(T_{k}\in\mathcal{C}\), so
\(\exc(T_{k})\le2\) by \eqref{eq:exc2eps}; for \(k\le3\) this is Lemma~\ref{lem:small}.

\emph{Reversal.} Put \(S_{k}:=T_{k}\) and, for \(2\le k\le N\), let \(x_{k}\) be the unique
cell of \(T_{k}\sm T_{k-1}\). Since \(T_{k}\) is connected with \(|T_{k}|\ge2\), the cell
\(x_{k}\) has a neighbour inside \(T_{k}\), necessarily in \(T_{k-1}=S_{k-1}\). This is (a),
and (b) was just shown.
\end{proof}

\begin{corollary}\label{cor:reach}
Let \(\mathcal{R}_{N}\) be the class of \(N\)-cell polyhexes reachable from a single cell by
a one-cell growth sequence all of whose intermediate configurations have edge-perimeter
excess at most \(2\), and \(\mathcal{M}_{N}\) the class of \(N\)-cell edge-perimeter
minimizers. Then \(\mathcal{M}_{N}\subseteq\mathcal{R}_{N}\) for every \(N\ge1\); that is,
restricting candidate generation to one-cell growth with excess at most \(2\) excludes no
minimizer. The reverse inclusion fails: \(\mathcal{R}_{N}\) contains many non-minimizers.
\end{corollary}


\section{Two obstructions}\label{sec:counter}

The proof above runs its induction inside an auxiliary class rather than inside the class of
minimizers, and pays for two devices that may look avoidable. Two explicit configurations
show that they are not.

\subsection{The corridor is a proper subset of the band}\label{sec:15.1}

\begin{example}\label{ex:9cell}
Let
\[
   X_{9}:=\{(0,1),(0,2),(0,3),(1,0),(1,1),(1,2),(1,3),(2,0),(2,2)\} .
\]
Then \(|X_{9}|=9\), \(\lambda(X_{9})=11=g(9)\) and \(R(X_{9})=12\), so
\(p(X_{9})=24=p_{\min}(9)+2\) and \(\exc(X_{9})=2\).
\end{example}

In words: the number of occupied lines is exactly the threshold, but one line carries two
runs instead of one, so the configuration is not directionally convex while its
edge-perimeter excess is still \(2\). All occupied lines have a single run except
\(q=2\), whose trace \(\{(2,0),(2,2)\}\) consists of two runs; thus
\(R(X_{9})=\lambda(X_{9})+1\).

Hence \(X_{9}\in\mathcal{B}_{2}\sm\mathcal{C}\), and the inclusion
\(\mathcal{C}\subsetneq\mathcal{B}_{2}\) of Section~\ref{sec:9.2} is strict. This is the
reason the corridor is defined with directional convexity built in rather than by the
excess bound alone: on \(\mathcal{B}_{2}\) the identity \(\exc=2\varepsilon\) is
unavailable, and the whole bookkeeping of the proof above rests on it.

\subsection{The corridor alone is not enough: a fifteen-cell trap}\label{sec:15.2}

\begin{example}\label{ex:15cell}
Let
\[
   X_{15}:=\{(0,4),(0,5)\}\cup\{(1,r):1\le r\le5\}\cup\{(2,r):0\le r\le4\}
   \cup\{(3,r):0\le r\le2\} .
\]
Then \(X_{15}\) is a directionally convex polyhex with \(|X_{15}|=15\),
\(\lambda(X_{15})=15\) and \(g(15)=14\), so \(\varepsilon(X_{15})=1\) and
\(X_{15}\in\mathcal{C}\) with \(\exc(X_{15})=2\). However
\[
   \rho(X_{15})=2,\qquad \tau(X_{15})=0,
\]
so \(\Adm(X_{15})\) fails, and every polyhex obtained from \(X_{15}\) by deleting one cell
has excess at least \(4\).
\end{example}

In words: the shortest extreme line still carries two cells while the plateau budget is
already exhausted, which is exactly the situation the admissibility condition forbids;
and at that point no single deletion can keep the excess at \(2\).

Consequently the naive statement

\begin{quote}
\emph{every polyhex of excess at most \(2\) with at least two cells has a cell whose
deletion leaves a polyhex of excess at most \(2\)}
\end{quote}

is false, and so is the version restricted to the corridor. The descent above
therefore cannot be run on \(\mathcal{C}\) alone: the pair
\(\mathcal{C}\cap\{\Adm\}\) is what closes, and Example~\ref{ex:15cell} is the witness that
the second condition is doing work. Note that this does not contradict
Proposition~\ref{pro:legal}, which asserts the existence of a connectedness-preserving
deletion; here such deletions exist, but none of them stays in the band.

\bibliographystyle{plain}
\bibliography{refs}
\end{document}
